\subsection{The Labors of Hercules, Part One}

Hercules\footnote{\url{https://en.wikipedia.org/wiki/Labours\_of\_Hercules}} is a picture of the strong man, who binds: in Christ's parables, he uses the image to refer to Satan, who always imitates God in a perverse manner. Esoterically, we might say that for most people (excluding Tomberg's mysterious ``just man''), ``Satan makes the first move in the chess game. A greater than Satan has to come to bind the strong man, and so Hercules is also a picture of ``the Lord'' or the ``true Self''. Since we live ``in Middle Earth'' \& find ourselves ``in media res'' (in the middle of things), there is work to do. Power comes, not necessarily all at once, but in a series of revelations. As the Western Tradition teaches, God ``respects'' free will -- we have to side with Him, or else the Asuras will bind us in service to Anti-Christ, under the dominion of a sunken Lucifer and an exalted Ahriman, the unholy Trinity.

This teaching even enters the popular imagination, in unexpected ways: in \emph{Star Wars}, the line is given -- ``No... there is another.'' This is Luke's ``sister'' (who, of course, seems more like a bride, at first), with whom the Force is strong. The Self has ``another Self'', \& also serves the ``Self Beyond the Self''.

Hercules is given tasks -- he has to attain heaven by storm. Yes, he is already Hercules, who as an infant, strangled the serpents in his crib. Yet, still, there is combat, storm, warfare. As Isaac Watts penned in the hymn, ``shall we look to go to heaven, on flowery beds of ease?''. However, it may be closer than we think, if we ride the storm.

His first task is to slay the Nemean Lion\footnote{\url{https://en.wikipedia.org/wiki/Nemean\_lion}}. King David had to slay a lion, it is worthy to note.

The Fathers mention, in the \emph{Philokalia}, several means of slaying demons. Some of the demons can be scorned or ridiculed, \& cannot bear focused attention on themselves, \& so the basic posture of ``awareness'', even letting the mind wander, to see where it goes, can bring fire upon their heads \& cause them to flee.

Which demon are we dealing with, here? The demon can change into a wounded damsel in distress, luring men into her cave, \& then devouring them. Since most men's sins revolve around sexual desire, we should naturally begin to examine this area of our lives? How \& why \& when are we tempted? In what manner?

The Biblical writers start the sequence with ``the lust of the flesh (and the pride of life)'' -- and this is the first temptation of the wilderness, to desire earthly food rather than God's Word. Women, for men, represent an appetite and a need, \& it is not too much to say that, for most men, they could not live without woman, and that much if not all of their Ego is bound up in psycho-sexual desire for the woman.

Even in popular culture, it is apparent\footnote{\url{http://alphagameplan.blogspot.com/2011/03/socio-sexual-hierarchy.html}} that the ``outsider'' status \& rough wanderings in the wilderness produces, not the Alpha Male, but the Sigma male, who is capable of attracting a beautiful consort, but in many senses, does not ``need'' or ``use'' her. This freedom from the woman, inner and outer, is a necessary prelude to true spiritual potency. Hence the medieval ideal of celibacy. If this appetite can be mastered, then one can engage the Nemean Lion. (Note, I am not endorsing, absolutely, these categories, but pointing out primarily that the rigid worship of ``Alpha'' status by males is misplaced, even as I acknowledge that their virtues are necessary, \& could be redeemed. It is they who ``set the tone'' of society, \& they are likely to be of a latent warrior class, who are ``masterless samurais'').

Scupoli advises the seeker to treat the lower chakra as a ``burning fire'' and to flee even the physical proximity of beautiful women, even women who are ``family'' or ``friends''. Since this is all but impossible today, the initiate must be even stronger. He must know, with absolute certitude, that the gravitational pull of beautiful women is a reflection of an inner illusion which must be absolutely put down, with all the ruthless of the Roman generals who sowed Carthage with salt.
\emph{Carthago delenda est}, was the founding of the ``Pax Romana''.


This does not disparage woman. Rather, it creates true freedom. What real man \& father would ever experience real temptation toward their own flesh and blood? None, if the man is healthy. Even so, the healthy initiate, in spirit, must view all woman, even the most desirable, as someone he must sweep the veil of sexual illusion from, in order to see properly, at all. Paradoxically, this will make him more more attractive to woman, \& invite further assault, for which he must be prepared, as a most cunning general. In ancient Christian tradition, the ``necessary lie'' was even advocated (eg., I am sick, I am mourning, etc.). Part of what drives the Kali Yuga is the misplaced sexual turmoil of alpha \& beta males (and to a lesser degree, everyone else's failure to not see through the illusion).

I would encourage all the men who read these articles to do whatever it takes to master the tiger of lust, absolutely. It will look differently, in various types and conditions. Chaste marriage to one woman is a form of celibacy, if exalted and properly nurtured. It absolutely constrains the sexual desire \emph{in all but one area, \&, eventually,} in modified form, even that area. The details \& circumstances will vary.

The mingling of races and flesh is inevitable in some ways, to some degrees, but in our time, it is actually the new Ideal. The initiate who comes to slay the Lion can take no part in this delusion, regardless of the consequences, personally. He has come to slay the Lion, not to sleep with the Demon. St. Benedict threw himself into thorns and brambles in order to master his lust. What have you done?

The power of the Holy Spirit that resides in the sexual chakra is typically the last center to be mastered by the initiate. It often remains, undisciplined, to the final breath. Along with pride, it is perhaps the most dangerous weakness of men (or women, for that matter). Man's creative organic power that resides there was meant to issue in the mouth in a word of power that commands and creates a new world (the mouth is the lower third of the head, and corresponds to the waist). It cannot do this if a man gives himself up to Lust, under the guise of ``riding the Tiger'', not perceiving that the meaning is twisted.

The foundation of asceticism is a denial of the need to quench our thirst for the female in anything other than God's sanction and God's time (which may be after our deaths, and not in the literal ``physical way''). It also cannot be accomplished without an invocation to Zeus (made explicit in the story of the Nemean lion, in which a sacrifice to Zeus is involved in the bargain). A higher desire is required to cast out the lower one, to make room for the action of God.

This is the foundation of power. This is Rome, who resisted ``Eastern luxury'', Phoenician materialism, \& the violent turbulence of the unprepared Northern tribes.. Rome, in the power of Ahriman (at that time, not overly exalted), held the line and upheld the true
\emph{virtus}. Hercules is a warrior and a priest, at once.


Issac the Syrian actually explains how this process of ``mortifying the flesh'' works:

\begin{quotex}
The effect of the cross is twofold; the duality of its nature divides it into two parts, \textbf{One consists in enduring sorrows of the flesh which are brought about by the action of the excitable part of the soul, and this part is called activity}. The \emph{\textbf{other part lies in the finer workings of the mind and in divine meditation, as well as in attending to prayer, etc.; it is accomplished by means of the desiring part of the soul and is called contemplation. The part of the soul by dint of its zeal, while the second part is the activity of soulful love, in other words, natural desire, which enlightens the rational part of the soul}}. \emph{Every man who, before perfectly mastering the first part, switches to the second, attracted out of weakness--to say nothing of laziness, is overtaken by God's wrath because he did not first mortify his members which are upon the earth (Col. 3:5).} In other words, he did not cure his thoughts of infirmities by patiently bearing the cross, but rather dared in his mind to envision the glory of the cross. 
\flright{\textit{Word} 55}
\end{quotex}


\begin{quotex}
It is evident from these words of Isaac the Syrian that what we call prelest proper exists when a man starts trying to live above his capabilities. Without having cleansed himself of passions, he strives for a life of contemplation and dreams of the delights of spiritual grace. Thus the wrath of God befalls a man; because he thinks too highly of himself, God's grace is withdrawn from him and he falls under the influence of the evil one who actively begins to tickle his vainglory with lofty contemplation and spiritual delights...\footnote{Source: \url{http://www.roca.org/OA/66-68/66n.htm}}
\end{quotex}

St Theophan offers even more detail:

\begin{quotex}
In the soul we find three powers: the intellect, the will, the heart, or, as the Holy Fathers say, the intellectual, desiring and incensive powers. Each of them is assigned particular curative exercises by the holy ascetics. These related excercises are both receptive and conducive to grace. They need not be contrived according to some theory, but rather chosen from tested ascetic labors particularly suited to a given power.
\flright{\textsc{St. Theophan the Recluse}, \textit{The Path To Salvation}}
\end{quotex}

Our wayward, random thoughts are ``women'' which attract the identification of Self: we chase them, they then transform into a demon, and ``slay'' us, by drawing us into illusions. It can be a long time before Hercules ever makes it out of the cave.


\flrightit{Posted on 2013-08-09 by Logres}

\begin{center}* * *\end{center}

\begin{footnotesize}
\begin{sffamily}
\texttt{August on 2013-08-10 at 07:31 said:}

A good article Logres.

Regarding the management of sexual energy, where it is not a question of controlled release through marriage or similar, there are diverse methods. Some involve exercises that open up the sexual centres, in a process that changes the normal polarity of sexual energy, turning it inward and harnessing its power to generate and sustain one's inner self, and this is not mere `sublimation' in the Freudian sense. Normally the energy seeks to flow outward and drive animal reproduction, or otherwise dissipate. Such processes can be rather dangerous for the untrained practitioner; consider how the ego is tossed around by normal levels of sexual stimulation, and imagine the strength required to resist very concentrated and strong charges of the same, directed inward. But these methods are very effective if mastered, they act even on the physical level.

Other ways involve blocking off or abandoning the lower centres, and this probably occurs in ascesis that begins on the psychological plane and tries to find a way beyond, ignoring and mortifying the flesh to silence its impulses. Sometimes the psychophysical substratum is ordered, sometimes merely restrained, with greater or lesser success.

A particular set of crises that might emerge on the path of sexual abstinence are the emotional complexes, particularly acute in the feminised men of today, many of whom are penetrated by obscure wounds that manifest as excessively sentimental and intimate tendencies, coupled with lack of inner form, which will begin to fester after the conscious mind decides to restrain its sexuality, even though, or perhaps because, the aspirant is affirming virility and strength through such an action. Most of this confusion originates in the mess of today's chaotic collective psychic and social atmosphere, which influences the formation of the individual. In any case, how one fares against these attacks is an important test, and the comments of Isaac the Syrian that you've quoted above are pertinent.

It must also be said that before the integration of the ego is accomplished, (good, quality) sex is important. It is one of the most potent activities that man can experience, and in those with proper capacity, it should bring to the surface elements of being that are deeper and more substantial than his unintegrated individuality, thus activating the need to be and prove oneself. This is the `rupture of levels' that Evola discusses in RTT, in this same context. It tends to occur somewhere beyond the crossfire of today's opposing sexual factions; those distorters who see nothing beyond the `animal ideal' and coarse physicality, and the sad reactions of that new caste of slovenly tattooed barbarian, the champion of ugliness.

The real man, who has laboured and attained more or less complete freedom from sexuality, may either set aside sex, or indulge as he sees fit.

\hfill

\texttt{Ash on 2013-08-11 at 03:32 said:}

The question is, what do we mean by ``freedom from sexuality''? Allowing for variance between the castes, Traditional teaching appears to dictate that this means being able to enjoy the legitimate union (say, a man and his lover) to its fullest extent without becoming enslaved or addicted to it. I read in a Muslim work on the matter of temptation, the title of which I forget, that it is considered as sinful to reject any pleasure given by God as it is to indulge what is haram. Evola writes in Revolt that the work of Tradition is to create channels for our chaotic currents to flow in the right direction, rather than the total elimination of those currents.

Writing from the Shia tradition, the Iranian Ali Shariati states it thus: ``However, with man being the battlefield, God and Satan are at war with each other. Thus, unlike former religions, the duality in Islam consists of worshiping two deities which exist in the constitution of man rather than in nature. Nature has a single deity and is under the dominance of only one God. This is why in Islam Satan is not standing against God but against the divine half of man. And since man is a two-dimensional creature who is kneaded of mud and God, he is in need of both. His ideology, religion, life, and civilization must all be capable of satisfying both of these dimensions. The tragedy is that history does not bear witness to this fact.''

Having struggled with this aspect of self-mastery, I would recommend that the tactic of laughing at the demon may be the one to pursue. The power of the sexual drive is difficult enough to master without us multiplying its strength by allowing it to weigh upon us like a milestone. Reflecting on childhood, it is far more effective to mock the ridiculous state of being controlled by it than to amplify its power by giving it the attraction of forbidden pleasure.

\hfill

\texttt{Cologero on 2013-08-11 at 10:18 said:}

August, regarding the ``real man'' that you mention, who presumably has transcended desire and is ruled by his intellect: under what circumstances, do you suppose, would he ``see fit'' to ``indulge''?

\hfill

\texttt{Gabriel on 2013-08-11 at 16:27 said:}

Step 15 of the Ladder of Divine Ascent by St. John Climacus is relevant.

``We have heard from that raving mistress gluttony, who has just spoken, that her offspring is war against bodily chastity. And this is not surprising, since our ancient forefather Adam teaches us this too. For if he had not been overcome by his stomach, he would not have known what a wife was (2). This is why those who keep the first commandment do not fall into the second transgression. And they continue to be children of Adam without knowing what Adam was. But they are made a little lower than the angels (in being subject to death.) And this is to prevent evil from becoming immortal.''

(2) i.e. he would have lived with her as with a sister.

``He is pure who expels love with love and who has extinguished the material fire by the immaterial fire.'' 15:2

``He who fights this adversary (lust) by bodily hardship and sweat is like one who has tied his foe with a string. But he who opposes him by temperance, sleeplessness, and vigil is like one who puts on a yoke on him. He who opposes him by humility, freedom from irritability, and thirst is like one who has killed his enemy and hidden him in the sand. And by sand, I mean humility, because it produces no fodder for the passions, but is mere earth and ashes.'' 15:12

``A fox pretends to be asleep, and the body and demons pretend to be chaste; the former in order to deceive a bird, and the latter in order to destroy a soul.'' 15:16

``He who has resolved to contend with his flesh and conquer it himself struggles in vain. For unless the Lord destroys the house of the flesh and builds the house of the soul, the person who wants to destroy it watches and fasts in vain.'' 15:25

``And our merciless foe, the teacher of fornication, says that our man-befriending God is very merciful towards this passion, since it is a natural one. But if we observe the guile of the demons, we shall find that, after sin has been committed, they say that God is a just and inexorable Judge.'' They said the former in order to lead us into sin, and now the latter to drown us in despair.''

``All demons try to darken our mind, and then they suggest what they like. For as long as the mind does not shut its eyes, we shall not be robbed of our treasure. But the demon of fornication tries to do this much more than all the rest. Often, after darkening out mind which governs us, it urges and disposes us in the presence of people to do what only those who are out of their mind do. Then later, when the mind becomes sober, we are ashamed of our unholy acts, words, and gestures, not only before those who saw us but also before ourselves, and we are amazed at our previous blindness. Often as a result of reflection, men has desisted this evil.'' 15:83

\hfill

\texttt{Gabriel on 2013-08-11 at 21:43 said:}

Since Adam and Eve married and had sexual relations after the fall, wouldn't the highest path still be to live as chaste, without marriage? Isn't this the path for the those ``not of this world''. Cologero wrote of this website ``We write for those who are both able and willing to abandon the contemporary world.'' I see the contemporary world as the world after the fall.

The Apostle Paul says in I Corinthians ``He that is unmarried careth for the things that belong to the Lord, how he may please the Lord: But he that is married careth for the things that are of the world, how he may please his wife. There is difference also between a wife and a virgin. The unmarried woman careth for the things of the Lord, that she may be holy both in body and in spirit: but she that is married careth for the things of the world, how she may please her husband'' (7:32-34). Marriage seems to be the lesser path, but acceptable. Earlier in the same chapter St. Paul writes ``For I would that all men were even as I myself. But every man hath his proper gift of God, one after this manner, and another after that. I say therefore to the unmarried and widows, It is good for them if they abide even as I. But if they cannot contain, let them marry: for it is better to marry than to burn'' (7:7-9).

Any ideas on what an ``exalted and properly nurtured'' chaste marriage would like? Fasting during fasts, raising children, not departing, etc? It seems like a fine line, how do you discern not having lust for your own wife, who is one in flesh with you? Even for those married or unmarried, how do you discern whether you chastely looking and conversing with a woman and lusting after a woman? How do you discern whether you are in this boat: ``Women, for men, represent an appetite and a need, \& it is not too much to say that, for most men, they could not live without woman, and that much if not all of their Ego is bound up in psycho-sexual desire for the woman.'' Could not even chastity be done out of pride for the sake of women?

One last point. Many Church Fathers bring up that if there would have been no fall, God would have populated the earth in a way different than sexual reproduction. Elder Thaddeus brings this up in Our Thoughts Determine Our Lives, Pg 126

``We live in families with out parents, brothers, and sisters and yet we are unsatisfied, we feel lonely and each one of us is burdened with his own cares. We want to leave our families and be with someone else; we long to cleave to that other person and spend our lives with him or her. God, when He created Adam, said that it is not good for man to be alone, and so He created a companion for him. Adam and Eve were one in the eyes of God.

But everything went wrong after the Fall. The Lord said ``Be fruitful and multiply, and replenish the earth (Gen 1:28). We do not know how reproduction would have taken place if man had not fallen. The earth was to have been filled with people (1).''

(1) According to St. Athanasius the Great (commentary on the Pslams 50:5), Gregory of Nyssa (On the Making of Man 17), John Chrysostom (On Virginity 14-15), Maximus Confessor (Ambiguum 41), John Damascene (Exact Exposition on the Orthodox Faith 4.24) and Symeon of Thessaloniki (On the Sacraments 38), if man had not fallen, God would have employed a means of increasing the human race other than sexual reproduction.

\hfill

\texttt{Logres on 2013-08-12 at 06:35 said:}

Gabriel, it can be a process, for most people: first of all, I agree that focusing on chastity can actually be a means of undermining it (which is why St Paul sanctioned marriage). ``It is better, etc., etc.''. If one lives with one's wife as if she was the ``polar being'', then (some have said), in some way, this will be given to you in eternity -- here is a possible means of avoiding the ``hell'' of Dante's lovers, but it runs by the edge of hell? It is clear, that in sex *usually'', the man ``loses to the woman''. If (however) the marriage is consecrated, \& the woman is submissive (as the KJV term is used), then, because of the ``one flesh'', the man has a chance to not lose spiritually -- therefore, it could be said that he loses, but it will be returned to him -- he will find himself again. I think this requires great love and holiness on both sides (although it might exist at first in a germinal form). Marriage, then, was a holy gamble. One way might be to realize it as a ``loss'', but to focus on the love that lays down its life -- in this way, at least, the male can de-identify with the lower urges, both by higher charity, \& marshalling of energy to sustain or undergo the ``loss''. The point would be to avoid the usual ``casual'' or pleasure-aspect of sex being uppermost, which is the default, and at least not to lie to one's self.

\hfill

\texttt{August on 2013-08-12 at 09:25 said:}

Cologero, I say it depends upon the man and the extent of his realisation.

Though not being subject to inherent desire, he might meet a woman whose privation seeks to make of him an axis, and if the conditions for deep attraction are sufficiently met, he could assume the role of male and consummate the bond. The woman would have brought desire into his presence, and he wills its assumption instead of letting it, and her, pass on. Such an act should not affect his deepest centre, which ought to remain steady if it is truly integral, and the dissolutive power of real women can be a great test of the extent to which this centre is established and capable of exercising control over peripheral elements of the man.

A properly qualified woman could also benefit in such a case, though I'm not sure she'd want to.

At the point where all relevant peripheral elements are totally disciplined and under spiritual control vis-a-vis sex, there is perhaps no longer anything for the man to gain from coupling; even the finest woman seems too coarse and clunky in comparison with the internal feminine. Except of course for social/family purposes, where appropriate.

\hfill

\texttt{Gabriel on 2013-08-12 at 12:05 said:}

Thank you for the answer Logres, that did bring clarity to me.

What you are proposing seems to be walking a very thin line August. I will stick with St. John Climacus for now.

``Let us, by every means in our power, avoid either seeing or hearing of that fruit which we have vowed not to taste. For it is absurd to think ourselves stronger than the Prophet David; that is impossible.'' 15:65 (2 Kgs. 11)

\hfill

\texttt{August on 2013-08-12 at 21:33 said:}

Yes, Gabriel, as thin as a razor's edge.

\hfill

\texttt{Logres on 2013-08-13 at 15:05 said:}

The Christian attitude to sex (at its best -- that is, in the meditation of those who truly know love/Love), has nothing to do with fear of the female or disgust at bodily functions : it has a great deal to do with distrust of thoughts ``as they appear'' to us, and are taken by the Ego. There are different methods or techniques, not to mention sacraments (eg., marriage), that have been articulated to sublimate and transform these thoughts. If Christianity has been uniquely (although not solely) responsible for sacralizing marriage, how is it that one can admit freely and without reserve the thought that it is repressive towards sex? Marriage was designed to affirm the naturalness of heterosexual relations, and also affirm chastity within that bond, not to mention providing for the attainment of ``polarity'', both through the intensity of the bond, and its ``eternal'' nature in this life. Indeed, in previous ages, people attained a species of moral exaltation in honoring vows which were unwisely taken. It was ``affirmation through negation'', which is a hallmark of Christian spirituality, generally. Don't confuse the corruption of the Ideal (indeed, what Ideal cannot be corrupted?), with its real perfection.

``I am not me.'' The narrow way is the broad path. That is the power of Hercules. I speak not here of other powers or other ways, which lead through this cave, as well, in the end. What is ``to master sex'' other than to achieve exactly what the sacrament of marriage achieves? I do not see the difficulty. Some will walk a different path, but (almost by definition), it cannot be the Ideal, except in a very special sense that does not translate well to other individuals. If I affirm that HO \& August are right, or can be in the path of the Spirit's wind (which I am inclined to do), does it follow that the lessons of the cave of Hercules (which the Desert Fathers set out to map \& master) are not the High Path? I am firmly convicted that the true romance of the path of Christ has been sung as of yet, in but bits and pieces : the mountain of the martyrs is surrounded by swords of paradise, and is crowned with a flame.

\hfill

\texttt{Logres on 2013-08-13 at 15:50 said:}

If you want to use Game-speak to phrase the task of counter-revolutionaries in our time, you might say that the Sigma male embodies some of the transformative possibilities available for young men in our time, in that the Sigma is capable of transforming himself and the world, rather than being conformed (which even, in our time) the Alpha most definitely is. The goal of a just order would be to ruled and guided by transformative Alphas, rather than degenerate ones. The stars turn.

\hfill

\texttt{Paulo on 2013-08-15 at 21:00 said:}

In the opinion of all of you,which type of expectations should we have toward women?I still don't have clear in my mind,what can i expect from them!And how about misogyny,are you all free from she?

\hfill

\texttt{Logres on 2013-08-16 at 10:44 said:}

You should expect, from one particular woman, that she will be your ``Other Self'' that symbolizes and makes real the ``Self Beyond the Self'' : whether she is literally so, or not (and this One does exist), if you treat her as such, she will become so virtually, \& will serve the purposes of reunion with this One. Other women, a handful, are compatible, or suitable, but are not this One, \& are meant more as sister-figures. The rest of women, you should be dead to, \& their allure is nothing more than biology masquerading as an illusion. Transcendent woman embodies the divine to man. Mouravieff calls this your ``polar being''. She is associated with the guardian angel that you have.

\hfill

\texttt{Paulo on 2013-08-16 at 13:13 said:}

Logres

And how about discovering the inner women who gonna complete me?Don't you think this is one possibility?So you really belive in''twin souls'',is that it?Is not polygamy the rule for us men,or this is just some modern western conditioning?And my final and most important question is,wicht type of psychological trait should i have,to be successful with women,dispassionate,gentile,fearful ect?Please answer me all these questions?

\hfill

\texttt{kid on 2013-08-16 at 17:06 said:}

man, logres that all sounds pretty mushy. you don't think marriage is best as a political or economic arrangement between estates or something, rather than some perfect union of two souls or whatever? i always imagined the latter idea of love to be something created in the last few hundred years where before people just got together with their neighbor at 15 and stuck with them forever because they had to.

\hfill

\texttt{Logres on 2013-08-16 at 21:41 said:}

Marriage can be karma yoga -- which is why Cologero spoke about the ``son of duty'' : you can't advance on your way, until you train (\& this often means fathering and raising) your replacement. So, yes, it's not necessary to dig too deeply into a practical marriage, and it can yield particularly bad results, when subverted (witness the modern ideal of ``romance'', pornographized by both males \& females, in their own way). Which of course produces a reaction, also very bad. I am speaking about what the Church aimed at in sacralizing it, not what it often happens or has to be; there are many ``practical'' marriages that were turned into romances by adherence to simple duty, over the years, which is a different story. We are assuming someone is already ``prepared'' to seek their other.

Paulo, it's difficult to give advice that is this personal. What do you mean by ``successful''? I would think benevolent detachment towards women is probably ideal, although this can be the WORST attitude if it falls into currying favor with females as some kind of religion. I think in older days, people were just more honest -- if you had to visit a house of ill-repute, you chalked it up to your own weakness, and went on. Today, men \& women play a big game, in which the Ego becomes attached. If a man wants a drink of water \& needs it, he takes it; the problem is, it is more complicated when the drink is another human being (St. Paul said ``they are one flesh''). It's not biological repulsion that guided the Church thinking on the subject, but a knowledge that connections were set up between individuals when they shared this together. Do you want to share karma with a total stranger you find hot? Many men do, \& you are also connected to them. So that is why detachment is useful -- how can you see what is happening if you buy into the game? BTW, not buying into the game (understood properly) seems to make you more attractive to women, not less. An unexpected side benefit, I suppose.

\hfill

\texttt{Paulo on 2013-08-16 at 22:46 said:}

What i mean by successful is to be THE MAN of the relationship,being able to control one self and to control the women in the most hard situations.What do you mean by''I would think benevolent detachment towards women is probably ideal, although this can be the WORST attitude if it falls into currying favor with females as some kind of religion.''?And''not buying into the game (understood properly).''?

\hfill

\texttt{Logres on 2013-08-16 at 23:22 said:}

Paulo, we shall continue on (your) thread in the Forum, so that we don't clutter the comments, and can have a more lengthy discussion... 

\hfill

\end{sffamily}
\end{footnotesize}

